\documentclass[10pt,a4paper]{extarticle}
\pagenumbering{gobble}
\usepackage[left=0.6in,top=0.6in,right=0.6in,bottom=0.6in]{geometry}
\usepackage[T1]{fontenc}
\usepackage[utf8]{inputenc}
\usepackage{helvet}
\renewcommand{\familydefault}{\sfdefault}
\usepackage[turkish,shorthands=off]{babel}
\usepackage[hidelinks]{hyperref}
\usepackage{titlesec}
\usepackage{enumitem}
\usepackage{datetime}
\usepackage{graphicx}

\titleformat{\section}{\large\bfseries\scshape\raggedright}{}{0em}{}[\titlerule]
\titlespacing*{\section}{0pt}{12pt}{8pt}

\newcommand{\entry}[2]{%
\noindent\begin{tabular*}{\textwidth}{@{\extracolsep{\fill}}lr}
\textbf{#1} & #2\\
\end{tabular*}
}
\newcommand{\subentry}[2]{%
\noindent\begin{tabular*}{\textwidth}{@{\extracolsep{\fill}}lr}
\textit{#1} & #2\\
\end{tabular*}
}

\begin{document}

\begin{center}
\begin{minipage}{\textwidth}
\centering
{\LARGE\textbf{Enes Orhan}} \hspace{2pt} {\LARGE{Backend \& Mobile Engineer}}\\[10pt]
\href{mailto:enesorhann404@gmail.com}{enesorhann404@gmail.com} \textbullet\
\href{https://linkedin.com/in/enesorhann}{linkedin.com/in/enesorhann} \textbullet\
\href{https://github.com/enesorhann}{github.com/enesorhann} \textbullet\
\href{https://enesorhan.com}{enesorhan.com}
\end{minipage}
\end{center}

\section{Özet}
14M+ ürün kaydı ve büyük ölçekli e-ticaret veri setleri üzerinde çalışan, ölçeklenebilir production sistemleri geliştirme deneyimine sahip Backend \& Mobile Engineer.
Yüksek performanslı API tasarımı (~200ms ortalama yanıt süresi), ML tabanlı microservice'ler ve Kubernetes tabanlı (kubeadm) deployment süreçlerinde uzmanlaştım.
Mimari tasarımdan performans optimizasyonuna, production hata ayıklamadan release yönetimine kadar backend yaşam döngüsünün tamamına sahiplik ettim.
10.000+ indirme sayısına ulaşan mobil uygulamalar geliştirdim ve binlerce kuruma hizmet veren B2B sistemlere destek verdim.

\section{Başarımlar}
\begin{itemize}[leftmargin=*,noitemsep,topsep=0pt]
\item 14M+ ürün kaydı üzerinde çalışan backend sistemlerini optimize ederek ortalama API yanıt süresini ~450ms'den ~200ms'ye düşürdüm.
\item 12M+ görsel, 529K+ kategori, 182K+ marka ve 3.344 mağaza içeren production sistemlerine katkı sağladım.
\item \%95 doğruluk oranına ulaşan ML tabanlı kategori eşleştirme microservice'i geliştirdim.
\item Redis caching ile FAISS tabanlı vektör arama pipeline'ı kurarak sorgulama gecikmesini iyileştirdim.
\item Süresi dolmuş premium kullanıcıların erişimini koruyabildiği abonelik doğrulama mantığındaki hatayı tespit ederek çözdüm.
\item kubeadm ile Kubernetes cluster'ları kurdum ve kubectl üzerinden pod hatalarını teşhis ettim.
\item 10.000+ toplam indirmeye ulaşan mobil uygulamalar yayınladım ve bakımını üstlendim.
\item Kendi sunucumuzda barındırılan mail altyapısını (SPF, DKIM, DMARC) tasarladım ve işlettim.
\end{itemize}

\section{Deneyim}

\entry{Freelance Software Developer}{Remote}
\subentry{Backend Engineer}{Ekim 2025 -- Devam ediyor}
\begin{itemize}[leftmargin=*,noitemsep,topsep=0pt]
\item RBAC tabanlı iş akışlarıyla 3.000+ şirketi destekleyen çok kiracılı (multi-tenant) Django REST backend mimarisi kurdum.
\item Kurum bazlı finansal operasyonları yöneten ~100 aktif kullanıcı için backend iş akışları tasarladım.
\item Param ve PayTR entegrasyonları için backend ödeme altyapısını hazırlayarak işlem doğrulama ve callback yönetim mantığını tasarladım.
\item Celery + Redis ile asenkron arka plan işleme süreçlerini hayata geçirdim.
\item Elasticsearch tabanlı filtreleme ve arama sistemlerini entegre ettim.
\item Tam e-posta kimlik doğrulama yığınıyla production seviyesinde mail sunucusu (Postfix + Dovecot) tasarladım ve işlettim.
\item kubeadm ile Kubernetes cluster'ları kurup yapılandırdım; deployment'ları ve production hata ayıklamayı kubectl üzerinden yönettim.
\item Mimari tasarımdan production release ve izlemeye kadar uçtan uca ticari backend sistemleri teslim ettim.
\end{itemize}

\entry{Fiyator A.Ş. \textit{(Ankara, Türkiye)} \textnormal{---} \href{https://fiyator.com}{fiyator.com}}{Remote}
\subentry{Full-Stack Mobile Developer Intern}{Haziran 2025 -- Ekim 2025}
\begin{itemize}[leftmargin=*,noitemsep,topsep=0pt]
\item 14M+ ürün kaydı, 12M+ görsel, 529K+ kategori ve 182K+ marka üzerinde çalışan backend sistemlerine katkı sağladım.
\item Çoğunluğu şirket hesaplarından oluşan 456 aktif kullanıcıya hizmet veren production ortamını destekledim.
\item PostgreSQL indeksleme ve Redis caching ile ortalama API yanıt süresini ~450ms'den ~200ms'ye düşürdüm.
\item Büyük ölçekli ürün veri setlerinde yüksek performanslı filtreleme sağlayan Elasticsearch entegrasyonunu gerçekleştirdim.
\item Semantic Transformers ve Jaccard benzerliği kullanarak \%95 doğruluğa ulaşan ML tabanlı kategori eşleştirme microservice'i (Django + FastAPI) geliştirdim.
\item Kategori sorgulama performansını iyileştirmek için Redis caching ile FAISS tabanlı vektör arama pipeline'ı tasarladım.
\item Akıllı ürün özelliği çıkarımı için Ollama (gemma3n:e2b) kullanan FastAPI microservice geliştirdim.
\item Süresi dolmuş premium kullanıcıların erişimini koruyabildiği abonelik doğrulama mantığındaki hatayı tespit ederek çözdüm.
\item Güvenli kimlik doğrulama ve token yenileme mekanizmaları (JWT interceptor mimarisi) tasarladım.
\item Kapsamlı testlerle 45+ fonksiyonel ve UI/UX sorununu tespit edip dokümante ettim.
\item Responsive UI iyileştirmelerini de kapsayan kapsamlı bir Flutter mobil yeniden tasarım sürecini yürüttüm.
\item Docker tabanlı deployment'lara ve production release döngülerine katıldım.
\end{itemize}

\section{Projeler}

\entry{BidsForShips (B2B Platform — Freelance) \textnormal{---} \href{https://bidsforships.com}{bidsforships.com}}{Ekim 2025 -- Şubat 2026}
\begin{itemize}[leftmargin=*,noitemsep,topsep=0pt]
\item RBAC tabanlı iş akışlarıyla 3.000+ şirketi destekleyen ölçeklenebilir Django REST backend geliştirdim.
\item ~100 aktif kullanıcı için backend iş akışları tasarladım.
\item Param ve PayTR entegrasyonları için backend ödeme altyapısını hazırladım.
\item Elasticsearch tabanlı arama ve filtreleme altyapısını entegre ettim.
\item Postfix + Dovecot kullanarak production seviyesinde mail sunucusu tasarladım ve işlettim.
\item Kubernetes ve CI/CD otomasyonu ile container tabanlı production yığınını deploy ettim.
\end{itemize}

\entry{Mail Sender Smart Edition \textnormal{---} \href{https://mailsendersmartedition.com}{mailsendersmartedition.com}}{Eylül 2025 -- Devam ediyor}
\begin{itemize}[leftmargin=*,noitemsep,topsep=0pt]
\item Şablon ve alıcı segmentasyonu mimarisine sahip ölçeklenebilir toplu e-posta yönetim platformu geliştirdim (20+ aktif kullanıcı).
\item Gemini 2.5 Flash ile LLM tabanlı e-posta şablonu üretimini hayata geçirerek optimize edilmiş token kullanımıyla yapılandırılmış JSON tabanlı taslaklar ürettim.
\item Zamanlama mantığına sahip, yeniden deneme tabanlı arka plan e-posta gönderim pipeline'ı tasarladım.
\item Güvenli ve güvenilir teslimat sağlayan, kendi sunucumuzda barındırılan mail altyapısını yapılandırdım.
\end{itemize}

\entry{Entegram (Mobil Uygulama)}{Kapalı Test}
\begin{itemize}[leftmargin=*,noitemsep,topsep=0pt]
\item Kullanıcıların günlük telefon alışkanlıklarını düzenli dil öğrenme oturumlarına dönüştürmek için tasarlanmış, kaydırma tabanlı bir İngilizce öğrenme uygulaması geliştirdim.
\item Sezgisel bir kaydırma etkileşimiyle kısa ve sık tekrarı teşvik eden, alışkanlık odaklı öğrenme akışı tasarladım.
\item Şu anda kapalı test sürecinde; yakın zamanda mağaza üzerinden genel kullanıma sunulması planlanıyor.
\end{itemize}

\section{Yetenekler}
\textbf{Diller:} Python, C\#, Dart, Kotlin, Java\\
\textbf{Backend:} Django, Django REST Framework, FastAPI, ASP.NET Core, Celery\\
\textbf{Mobil:} Flutter, Firebase\\
\textbf{Veritabanları:} PostgreSQL, SQLite, Elasticsearch, Redis\\
\textbf{DevOps:} Docker, Kubernetes (kubeadm), GitHub Actions, Linux, CI/CD, Production Debugging\\
\textbf{Kavramlar:} REST Architecture, RBAC, Veri Modelleme, Caching Stratejileri, ML Pipelines

\section{Eğitim}
\textbf{Konya Teknik Üniversitesi} \hfill Konya, Türkiye\\
\textit{Bilgisayar Mühendisliği, Lisans} \hfill Ağustos 2022 -- Haziran 2026\\
Not Ortalaması: 3.40 / 4.00

\end{document}
